\documentclass[11pt, a4paper]{article}
\usepackage[a4paper, margin=2.5cm]{geometry}
\usepackage{fontspec}
\usepackage[german, bidi=basic, provide=*]{babel}
\babelfont{rm}{Noto Sans}
\usepackage{booktabs}
\usepackage{xcolor}
\usepackage{enumitem}
\usepackage{fancyhdr}

\definecolor{bhtblue}{RGB}{0, 51, 102}
\definecolor{bhtgray}{RGB}{100, 116, 139}

\pagestyle{fancy}
\fancyhf{}
\fancyhead[L]{\small\color{bhtgray}Projektantrag Laboreinrichtung | F.-o. E-Projekt Mechatronik}
\fancyhead[R]{\small\color{bhtgray}BHT Berlin}
\fancyfoot[C]{\thepage}
\renewcommand{\headrulewidth}{0.4pt}
\renewcommand{\footrulewidth}{0pt}

\begin{document}

\begin{center}
    {\LARGE \textbf{\textcolor{bhtblue}{Detaillierte Beschreibung der Arbeitspakete}}}\\[6pt]
    {\large \textbf{Entwicklung eines Kraft-Weg-Messplatzes für Kraftsensoren aus leitendem PLA als DMS}}\\[4pt]
    \normalsize \textit{Betreuer: Prof. Dr.-Ing. Karsten Pietsch | Labor B402L (GOS)}
\end{center}

\vspace{0.5em}

\noindent Dieses Dokument beschreibt die strukturierten Arbeitspakete (AP) für das studentische Mechatronik-Projekt (2--3 Studierende, Budget $<200\,€$) an der Berliner Hochschule für Technik.

\section*{Übersicht der Projektphasen und Arbeitspakete}

\subsection*{AP 1: Projektmanagement, Konzeption und Anforderungsanalyse}
\begin{itemize}[leftmargin=*, label=$\bullet$]
    \item \textbf{Ziele:} Definition der genauen Systemanforderungen, Zeit- und Ressourcenplanung sowie strikte Budgetüberwachung ($<200\,€$).
    \item \textbf{Inhalte:}
    \begin{itemize}[leftmargin=*]
        \item Erstellung eines Lasten- und Pflichtenhefts für den Messplatz.
        \item Recherche zum Stand der Technik (insb. zu additiv gefertigten Kraftsensoren aus leitendem PLA, basierend auf der Vorlage von Elgeneidy et al., 2018).
        \item Definition der mechanischen, elektrischen und softwareseitigen Schnittstellen.
    \end{itemize}
\end{itemize}

\subsection*{AP 2: Mechanische Konstruktion und CAD-Entwicklung (CAD)}
\begin{itemize}[leftmargin=*, label=$\bullet$]
    \item \textbf{Ziele:} Entwurf, Konstruktion und physische Realisierung des mechanischen Prüfstands zur präzisen Einspannung und Belastung der Sensoren.
    \item \textbf{Inhalte:}
    \begin{itemize}[leftmargin=*]
        \item CAD-Modellierung der Führungseinrichtung und der Probenhalterungen.
        \item Konstruktion der mechanischen Schnittstellen für den Kraftaufnehmer und Verstellelemente.
        \item Auswahl und Beschaffung der mechanischen Standardkomponenten innerhalb des Budgets.
        \item Fertigung bzw. 3D-Druck der spezifischen Prototypenteile.
    \end{itemize}
\end{itemize}

\subsection*{AP 3: Elektrotechnik und Sensoranbindung (ET)}
\begin{itemize}[leftmargin=*, label=$\bullet$]
    \item \textbf{Ziele:} Schaltungstechnische Integration der additiv gefertigten Sensoren und Signalaufbereitung.
    \item \textbf{Inhalte:}
    \begin{itemize}[leftmargin=*]
        \item Sichere Kontaktierung der Kraftsensoren aus leitendem PLA (mittels leitfähigem Faden und Metallstiften).
        \item Aufbau der Messschaltung (z.\,B. Wheatstone-Brücke zur Erfassung kleinster Widerstandsänderungen bei Verformung).
        \item Verkabelung, Signalverstärkung und Anbindung an einen Mikrocontroller (z.\,B. Arduino oder ESP32).
    \end{itemize}
\end{itemize}

\subsection*{AP 4: Softwareentwicklung und Steuerung (SW)}
\begin{itemize}[leftmargin=*, label=$\bullet$]
    \item \textbf{Ziele:} Programmierung der Steuerungs-, Mess- und Auswertesoftware zur Erfassung der physikalischen Größen.
    \item \textbf{Inhalte:}
    \begin{itemize}[leftmargin=*]
        \item Entwicklung der Firmware für die Datenerfassung am Mikrocontroller.
        \item Erstellung einer grafischen Benutzeroberfläche oder eines Skripts (z.\,B. Python oder MATLAB) zur Echtzeit-Visualisierung der Kraft-Weg-Kurven.
        \item Implementierung von Filter- und Kalibrierungsalgorithmen.
    \end{itemize}
\end{itemize}

\subsection*{AP 5: Gesamtsystemintegration, Kalibrierung und Dokumentation}
\begin{itemize}[leftmargin=*, label=$\bullet$]
    \item \textbf{Ziele:} Zusammenführung aller Teilsysteme, Validierung der Messgenauigkeit und Bereitstellung als funktionierender Labor-Demonstrator.
    \item \textbf{Inhalte:}
    \begin{itemize}[leftmargin=*]
        \item Hardware-in-the-Loop-Tests des Gesamtaufbaus.
        \item Kalibrierung des Messplatzes mit definierten Referenzlasten und -wegen.
        \item Erstellung einer umfassenden Projektdokumentation und einer Bedienungsanleitung für kommende Studierendengenerationen.
        \item Abschlusspräsentation und offizielle Übergabe des Demonstrators an das Labor.
    \end{itemize}
\end{itemize}

\vspace{1em}
\hrule
\vspace{0.5em}
\noindent \textit{Hinweis: Dieses Dokument dient als strukturierte Arbeitsgrundlage für den Projektantrag und wird hochschulintern verwendet.}

\end{document}